\documentclass{beamer}

% Configuración del tema
\usetheme{Madrid}
\usecolortheme{default}

% Paquetes para español y codificación
\usepackage[utf8]{inputenc}
\usepackage[T1]{fontenc}
\usepackage[spanish]{babel}
\usepackage{graphicx}
\usepackage{tikz}
\usetikzlibrary{positioning, arrows.meta, shapes.geometric}
\usepackage{hyperref}

% Información del título
\title[Clase 2: Requerimientos]{Clase 2: Análisis de Requerimientos Funcionales y No Funcionales}
\subtitle{Proyecto Web}
\author{Rodrigo García, PhD}
\date{\today}

\begin{document}

% Diapositiva de Título
\begin{frame}
    \titlepage
\end{frame}

% Índice
\begin{frame}{Agenda de la Clase}
    \tableofcontents
\end{frame}

% ==========================================
% SECCIÓN 1: INTRODUCCIÓN
% ==========================================
\section{Introducción a los Requerimientos}

\begin{frame}{¿Qué son los Requerimientos?}
    \begin{itemize}
        \item Un requerimiento es una condición o capacidad que debe cumplir un sistema para satisfacer un contrato, estándar o especificación.
        \item Son la base de cualquier proyecto de desarrollo de software.
        \item \textbf{Propósito:} Asegurar que lo que construimos es exactamente lo que el cliente o el usuario final necesita.
    \end{itemize}
\end{frame}

\begin{frame}{Importancia de unos Buenos Requerimientos}
    \begin{itemize}
        \item \textbf{Reducen costos y tiempo:} Detectar un error en la fase de requerimientos es mucho más barato que corregirlo en producción.
        \item \textbf{Guían el diseño:} Ayudan a definir la arquitectura del sistema, la base de datos y la interfaz de usuario.
        \item \textbf{Miden el éxito:} Sirven como criterio de aceptación para validar que el proyecto fue exitoso.
    \end{itemize}
\end{frame}

% ==========================================
% SECCIÓN 2: REQUERIMIENTOS FUNCIONALES
% ==========================================
\section{Requerimientos Funcionales (RF)}

\begin{frame}{¿Qué son los Requerimientos Funcionales?}
    \begin{itemize}
        \item Definen \textbf{LO QUE} el sistema debe hacer.
        \item Son los comportamientos, funciones y procesos específicos que el software debe ser capaz de ejecutar.
        \item Describen la interacción entre el sistema y su entorno (usuarios, otros sistemas).
    \end{itemize}
    \vspace{0.3cm}
    \begin{block}{Regla mnemotécnica}
    \textbf{Verbos:} El sistema \textbf{debe permitir}..., El sistema \textbf{calculará}..., El usuario \textbf{podrá iniciar sesión}....
    \end{block}
\end{frame}

\begin{frame}{Ejemplos de Requerimientos Funcionales}
    \begin{itemize}
        \item El sistema debe permitir a los usuarios registrarse con un correo electrónico y una contraseña.
        \item El sistema debe enviar un correo de confirmación al usuario tras un registro exitoso.
        \item El administrador podrá generar un reporte mensual de ventas en formato PDF.
        \item El sistema debe permitir al usuario añadir productos a un carrito de compras.
        \item El software calculará automáticamente el IVA aplicable a la compra final.
    \end{itemize}
\end{frame}

% ==========================================
% SECCIÓN 3: REQUERIMIENTOS NO FUNCIONALES
% ==========================================
\section{Requerimientos No Funcionales (RNF)}

\begin{frame}{¿Qué son los Requerimientos No Funcionales?}
    \begin{itemize}
        \item Definen \textbf{CÓMO} el sistema debe comportarse o las cualidades que debe poseer.
        \item No hablan de una acción específica, sino de las \textbf{restricciones} y \textbf{atributos de calidad} del sistema.
    \end{itemize}
    \vspace{0.3cm}
    \begin{block}{Categorías comunes}
    Rendimiento, Seguridad, Usabilidad, Escalabilidad, Disponibilidad, Compatibilidad.
    \end{block}
\end{frame}

\begin{frame}{Ejemplos de Requerimientos No Funcionales}
    \begin{itemize}
        \item \textbf{Rendimiento:} La página de inicio debe cargar en menos de 2 segundos.
        \item \textbf{Seguridad:} Las contraseñas de los usuarios deben almacenarse encriptadas en la base de datos (Ej: bcrypt).
        \item \textbf{Usabilidad:} El sistema debe ser responsivo y verse bien tanto en móviles como en computadoras de escritorio.
        \item \textbf{Disponibilidad:} El sistema debe estar disponible el 99.9\% del tiempo.
        \item \textbf{Compatibilidad:} La aplicación web debe ser compatible con las últimas versiones de Chrome, Firefox y Safari.
    \end{itemize}
\end{frame}

% ==========================================
% SECCIÓN 4: DIFERENCIAS Y EJERCICIOS
% ==========================================
\section{Diferencias Clave y Ejercicios}

\begin{frame}{Diferencia Principal (Resumen)}
    \begin{center}
    \begin{tabular}{|p{5cm}|p{5cm}|}
    \hline
    \textbf{Funcionales (RF)} & \textbf{No Funcionales (RNF)} \\
    \hline
    Qué debe hacer el sistema. & Cómo debe comportarse el sistema. \\
    \hline
    Si no se cumple, el sistema no funciona. & Si no se cumple, el sistema puede funcionar pero con mala calidad. \\
    \hline
    Son específicos de la aplicación. & Suelen ser generales para la arquitectura. \\
    \hline
    \end{tabular}
    \end{center}
\end{frame}

\begin{frame}{Próximos Pasos - Actividad del Proyecto}
    \begin{block}{Tarea Práctica}
        \begin{enumerate}
            \item Reúnete con tu equipo y tomen la idea del proyecto definida en la Clase 1.
            \item Redacten al menos \textbf{5 Requerimientos Funcionales} principales de su aplicación.
            \item Redacten al menos \textbf{5 Requerimientos No Funcionales} que consideren críticos.
        \end{enumerate}
    \end{block}
    \vspace{0.5cm}
    \begin{center}
        \Large \textbf{¡Sigamos construyendo!}
    \end{center}
\end{frame}

\end{document}
